% ============================================================
%  Глава 6: Заключение
% ============================================================

\section*{Заключение}
\addcontentsline{toc}{section}{Заключение}
\label{sec:conclusion}

\subsection*{Основные результаты работы}

В данной магистерской диссертации разработана и исследована система
\textbf{быстрой калибровки модели грубой волатильности Rough Heston}
на основе операторных нейронных сетей (FNO) с FiLM-кондиционированием
и автоградиентного метода Ньютона--Гаусса. Все поставленные задачи решены.

\medskip
\textbf{1. GPU-ускоренный прайсер Фурье--COS.}
Реализован устойчивый численный решатель системы поднятых ОДУ Rough Heston
с использованием экспоненциального полушагового интегратора (exponential
midpoint rule). Решатель безусловно устойчив по линейному члену, обрабатывает
батч из 1024 наборов параметров за $\approx 1{,}4\,\text{с}$ на RTX~3060.
Точность при $N=40$ факторах Бернштейна: максимальная ошибка IV $37{,}9\,\text{б.п.}$,
при $N_{\text{steps}}=200$ --- погрешность ODE $1{,}9\,\text{б.п.}$ при $T=0{,}1$.

\medskip
\textbf{2. Суррогатная модель FNO v2.}
Обучена суррогатная модель \texttt{MirrorPaddedFNO2d} с FiLM-кондиционированием
на 50\,000 образцах Sobol. Итоговые метрики:
\[
  R^2 = 0{,}9991, \quad \text{MAE} = 0{,}058\%, \quad
  \text{инференс} < 1\,\text{мс} \quad (1400\times\text{ быстрее прямого прайсера}).
\]

\medskip
\textbf{3. Анализ идентифицируемости (FIM).}
Проведён численный анализ матрицы Фишера для модели Rough Heston.
Переход от 6D к 3D-репараметризации $(v_0, \zeta, \lambda)$ снижает
число обусловленности задачи в $\mathbf{1301\times}$~---
с $\kappa \sim 10^6$ до $\kappa \approx 770$~--- теоретически обосновывая
фиксацию призрачных параметров $(\kappa, \theta, H)$.

\medskip
\textbf{4. Метод Ньютона--Гаусса с \texttt{jacfwd}.}
Разработан быстрый калибратор на основе метода Ньютона--Гаусса
с якобианом, вычисленным через прямое автодифференцирование
(\texttt{torch.func.jacfwd}). По сравнению с L-BFGS:
\begin{itemize}
  \item в $3\times$ быстрее ($541$ vs $1599\,\text{мс}$),
  \item точнее по $v_0$: $1{,}7\%$ vs $4{,}3\%$ (при нулевом шуме),
  \item более робастен: сходимость $96{,}7\%$ при всех уровнях шума.
\end{itemize}

\medskip
\textbf{5. Реал-тайм потоковая калибровка.}
Продемонстрирована потоковая калибровка на 20 синтетических тиках SPX-опционов
с шумом $1\%$: $p_{95} = 668\,\text{мс} < 1000\,\text{мс}$~--- система
обрабатывает $95\%$ тиков в реальном времени с запасом $1{,}5\times$
относительно частоты обновления котировок ($\approx 1$--$2\,\text{с}$).

\medskip
\textbf{6. Дельта-хеджирование.}
Бэктест на 30 путях Монте-Карло (горизонт 1~год, 52 шага): FNO-$\Delta$
снижает дисперсию P\&L на $5{,}1\%$ относительно плоской BS-дельты,
при среднем времени калибровки $518\,\text{мс/шаг}$.

\medskip
\textbf{7. Обучаемый показатель Хёрста (FNO v3).}
Расширение пространства параметров до 6D с включением $H \in [0{,}04;\,0{,}15]$
реализовано, обучено (500 эпох, SWA, 30{,}4 мин) и оценено:
\[
  R^2_{\text{genuine}} = 0{,}9981, \quad
  \text{MAE}_{\text{genuine}} = 0{,}264\%, \quad
  \text{MAE}_{\text{all}} = 0{,}333\%.
\]
Точность в $4{,}5\times$ ниже v2 из-за расширения пространства параметров
и $71\%$ медианных меток при $T=0{,}1$~--- однако MAE $< 0{,}3\%$ ($30\,\text{б.п.}$)
остаётся ниже рыночного bid-ask спреда и достаточна для практической калибровки.
Реализован 4D-калибратор $(v_0, \zeta, \lambda, H)$ с автоадаптацией
показателя Хёрста к текущему режиму рынка.

\subsection*{Научная новизна подтверждена}

\begin{enumerate}
  \item FiLM-FNO для прайсинга Rough Heston~--- первое применение FNO с
        FiLM-кондиционированием для IV-поверхностей с дробным броуновским движением.
  \item Автоградиентный GN-калибратор~--- применение \texttt{jacfwd} через
        нейронный оператор с квадратичной сходимостью.
  \item FIM-обоснование редукции параметров: $1301\times$ снижение числа
        обусловленности.
  \item Обучаемый $H$: расширение суррогата для переменного показателя Хёрста.
\end{enumerate}

\subsection*{Ограничения и направления дальнейших исследований}

\begin{itemize}
  \item \textbf{Нестационарность}: фиксация $(\kappa^*, \theta^*, H^*)$
        вносит смещение при значительном отклонении реального рынка от этих значений.
        Направление: адаптивное обновление фиксированных параметров по долгосрочным
        опционам (LEAPS) раз в месяц.

  \item \textbf{Экзотические деривативы}: суррогат обучен только на
        европейских ванильных опционах (сетка $8\times11$). Расширение
        на барьерные, азиатские и другие экзотики требует новых датасетов
        и, возможно, 3D-FNO (добавляется ось базового актива).

  \item \textbf{Масштаб}: RTX~3060 ограничивает батч $B\leq1024$.
        На серверных GPU (A100, H100) ожидается $10$--$50\times$ ускорение,
        что откроет возможность онлайн-калибровки всей книги опционов.

  \item \textbf{Неопределённость}: методы Bayesian Neural Operator
        (Magnani~et~al., 2022)~\cite{Magnani2022} позволяют получить
        доверительные интервалы для IV-поверхности~--- важный аспект для
        риск-менеджмента.
\end{itemize}

\subsection*{Заключительное слово}

Предложенная система демонстрирует, что \textbf{нейронные сети могут
кардинально ускорить калибровку сложных моделей грубой волатильности}
без существенной потери точности~--- открывая путь к применению
экономически более корректных моделей в реал-тайм трейдинге и риск-менеджменте.
